\chapter{Template Style Guide}
\label{chap:style_guide}

This chapter serves as a visual guide to the capabilities of the Thesis Template. It demonstrates the default styling for tables, mathematics, code listings, and image grids, ensuring all characteristics are properly showcased.

\section{Advanced Tables}
Tables should be clean and professional, avoiding vertical lines where possible. The \texttt{booktabs} package provides high-quality horizontal rules.

\begin{table}[H]
    \centering
    \caption{Component Selection Matrix (Booktabs Style)}
    \label{tab:selection_matrix_showcase}
    \begin{tabularx}{\textwidth}{@{} l X c c @{}}
        \toprule
        \textbf{Component} & \textbf{Description} & \textbf{Cost (\$)} & \textbf{Selected} \\
        \midrule
        Microcontroller & ESP32 Dual Core, WiFi/BT & 5.00 & Yes \\
        Motor Driver & L298N Dual H-Bridge & 2.50 & No \\
        Motor Driver & TB6612FNG High Efficiency & 3.00 & Yes \\
        Battery & LiPo 11.1V 2200mAh & 15.00 & Yes \\
        \bottomrule
    \end{tabularx}
\end{table}

\begin{table}[H]
    \centering
    \caption{Complex Multi-column Table}
    \label{tab:multicol_showcase}
    \begin{tabular}{l c c c}
        \toprule
        & \multicolumn{3}{c}{\textbf{Performance Metrics}} \\
        \cmidrule(lr){2-4}
        \textbf{Algorithm} & \textbf{Accuracy (\%)} & \textbf{Latency (ms)} & \textbf{CPU Load (\%)} \\
        \midrule
        A* Search & 99.5 & 120 & 15 \\
        Dijkstra & 100.0 & 450 & 45 \\
        RRT* & 95.0 & 25 & 8 \\
        \bottomrule
    \end{tabular}
\end{table}

\section{Mathematics and Equations}
The template uses \texttt{newtxmath} for professional mathematical typesetting.

\subsection{Standard Equations}
Maxwell's equations in differential form:
\begin{align}
    \nabla \cdot \mathbf{E} &= \frac{\rho}{\varepsilon_0} \label{eq:maxwell1} \\
    \nabla \cdot \mathbf{B} &= 0 \label{eq:maxwell2} \\
    \nabla \times \mathbf{E} &= -\frac{\partial \mathbf{B}}{\partial t} \label{eq:maxwell3} \\
    \nabla \times \mathbf{B} &= \mu_0\mathbf{J} + \mu_0\varepsilon_0\frac{\partial \mathbf{E}}{\partial t} \label{eq:maxwell4}
\end{align}

\subsection{Matrices}
State-space representation of the robot dynamics:
\begin{equation}
    \begin{bmatrix}
        \dot{x} \\
        \dot{y} \\
        \dot{\theta}
    \end{bmatrix}
    =
    \begin{bmatrix}
        \cos\theta & 0 \\
        \sin\theta & 0 \\
        0 & 1
    \end{bmatrix}
    \begin{bmatrix}
        v \\
        \omega
    \end{bmatrix}
\end{equation}

\subsection{Theorems and Definitions}
\begin{notebox}
    \textbf{Definition (Swarm Intelligence):} The collective behavior of decentralized, self-organized systems, natural or artificial.
\end{notebox}

\section{Code Listings}
Source code is formatted using the \texttt{listings} package with syntax highlighting.

\begin{lstlisting}[language=Python, caption={Python Control Loop Example}, label={lst:python_example}]
import numpy as np

def control_loop(target_pose, current_pose):
    """
    Computes velocity commands to reach target pose.
    """
    error = target_pose - current_pose
    kp_linear = 0.5
    kp_angular = 2.0
    
    v = kp_linear * error[0]
    omega = kp_angular * error[2]
    
    return v, omega
\end{lstlisting}

\section{Image Grids and Layouts}
We demonstrate various arrangements of figures using the \texttt{subcaption} package standard.

\subsection{1x2 Grid (Side-by-Side)}
\begin{figure}[H]
    \centering
    \begin{subfigure}[b]{0.48\textwidth}
        \centering
        \includegraphics[width=\textwidth]{example-image-a}
        \caption{Left View}
        \label{fig:side_a}
    \end{subfigure}
    \hfill
    \begin{subfigure}[b]{0.48\textwidth}
        \centering
        \includegraphics[width=\textwidth]{example-image-b}
        \caption{Right View}
        \label{fig:side_b}
    \end{subfigure}
    \caption{Two images displayed side-by-side.}
    \label{fig:1x2_grid}
\end{figure}

\subsection{2x2 Grid}
\begin{figure}[H]
    \centering
    \begin{subfigure}[b]{0.45\textwidth}
        \centering
        \includegraphics[width=\textwidth]{example-image-a}
        \caption{Top Left}
    \end{subfigure}
    \hfill
    \begin{subfigure}[b]{0.45\textwidth}
        \centering
        \includegraphics[width=\textwidth]{example-image-b}
        \caption{Top Right}
    \end{subfigure}
    \vskip\baselineskip
    \begin{subfigure}[b]{0.45\textwidth}
        \centering
        \includegraphics[width=\textwidth]{example-image-c}
        \caption{Bottom Left}
    \end{subfigure}
    \hfill
    \begin{subfigure}[b]{0.45\textwidth}
        \centering
        \includegraphics[width=\textwidth]{example-image}
        \caption{Bottom Right}
    \end{subfigure}
    \caption{A 2x2 grid of four images.}
    \label{fig:2x2_grid}
\end{figure}

\subsection{4x4 Complex Grid (Small Multiples)}
Evaluating results across many scenarios often requires denser grids.

\begin{figure}[H]
    \centering
    % Row 1
    \begin{subfigure}[b]{0.23\textwidth}\includegraphics[width=\textwidth]{example-image-a}\caption*{1}\end{subfigure}\hfill
    \begin{subfigure}[b]{0.23\textwidth}\includegraphics[width=\textwidth]{example-image-b}\caption*{2}\end{subfigure}\hfill
    \begin{subfigure}[b]{0.23\textwidth}\includegraphics[width=\textwidth]{example-image-c}\caption*{3}\end{subfigure}\hfill
    \begin{subfigure}[b]{0.23\textwidth}\includegraphics[width=\textwidth]{example-image}\caption*{4}\end{subfigure}
    \vskip+1ex
    % Row 2
    \begin{subfigure}[b]{0.23\textwidth}\includegraphics[width=\textwidth]{example-image}\caption*{5}\end{subfigure}\hfill
    \begin{subfigure}[b]{0.23\textwidth}\includegraphics[width=\textwidth]{example-image-a}\caption*{6}\end{subfigure}\hfill
    \begin{subfigure}[b]{0.23\textwidth}\includegraphics[width=\textwidth]{example-image-b}\caption*{7}\end{subfigure}\hfill
    \begin{subfigure}[b]{0.23\textwidth}\includegraphics[width=\textwidth]{example-image-c}\caption*{8}\end{subfigure}
    \vskip+1ex
    % Row 3
    \begin{subfigure}[b]{0.23\textwidth}\includegraphics[width=\textwidth]{example-image-c}\caption*{9}\end{subfigure}\hfill
    \begin{subfigure}[b]{0.23\textwidth}\includegraphics[width=\textwidth]{example-image}\caption*{10}\end{subfigure}\hfill
    \begin{subfigure}[b]{0.23\textwidth}\includegraphics[width=\textwidth]{example-image-a}\caption*{11}\end{subfigure}\hfill
    \begin{subfigure}[b]{0.23\textwidth}\includegraphics[width=\textwidth]{example-image-b}\caption*{12}\end{subfigure}
    
    \caption{A dense grid layout for comparing multiple data points or results.}
    \label{fig:dense_grid}
\end{figure}
